\documentclass[12pt,letterpaper]{article}
\usepackage[utf8]{inputenc}
\usepackage[spanish]{babel}
\usepackage{geometry}
\usepackage{fancyhdr}
\usepackage{graphicx}
\usepackage{xcolor}
\usepackage{amsmath}
\usepackage{parskip}
\usepackage{titlesec}
\usepackage{setspace}
\usepackage{float}
\usepackage{newtxtext,newtxmath}
\usepackage{array}
\usepackage{longtable}
\usepackage{hyperref}

% Configuración de página
\geometry{
    top=2.5cm,
    bottom=2.5cm,
    left=2.5cm,
    right=2.5cm,
    headheight=1.5cm,
    footskip=1cm
}

% Configuración de títulos
\titleformat{\section}{\large\bfseries}{\thesection.}{1em}{}
\titleformat{\subsection}{\normalsize\bfseries}{\thesubsection)}{1em}{}

% Configuración de encabezados y pies de página
\pagestyle{fancy}
\fancyhf{}
\fancyhead[L]{\small \textit{\nombreproyecto}}
\fancyhead[R]{\small \materia}
\fancyfoot[C]{\thepage}
\renewcommand{\headrulewidth}{0.4pt}

% =====================================================
% VARIABLES - PERSONALIZAR AQUÍ
% =====================================================
\newcommand{\universidad}{UNIVERSIDAD AUTÓNOMA DE SAN LUIS POTOSÍ}
\newcommand{\facultad}{Facultad de Ingeniería}
\newcommand{\carrera}{Ingeniería en Computación}
\newcommand{\materia}{Aplicaciones Web Escalables (224501)}
\newcommand{\nombreproyecto}{KitchenFlow: Control de Producción y Costeo para Cafeterías}
\newcommand{\autores}{Miguel Alejandro Gutiérrez Silva}
\newcommand{\fecha}{Semestre 2025-2026/II}
\newcommand{\profesor}{Ing. Francisco Javier Gómez Vázquez}

\begin{document}

\onecolumn
\singlespacing

% =====================================================
% PÁGINA DE PRESENTACIÓN
% =====================================================
\thispagestyle{empty}

\begin{center}
	\vspace*{0.5cm}

	{\color{black!70}\rule{\textwidth}{1.5pt}}

	\vspace{0.8cm}

	{\LARGE \textbf{\universidad}}

	\vspace{0.5cm}

	{\Large \facultad}

	\vspace{0.3cm}

	{\large \carrera}

	\vspace{0.8cm}

	{\color{black!40}\rule{0.6\textwidth}{0.5pt}}

	\vspace{0.8cm}

	{\large Propuesta de Proyecto}

	\vspace{0.2cm}

	{\Large \textbf{\materia}}

	\vspace{1.5cm}

	\fbox{%
		\parbox{0.85\textwidth}{%
			\vspace{0.5cm}
			\centering
			{\LARGE \textbf{\nombreproyecto}}
			\vspace{0.5cm}
		}%
	}

	\vspace{1.5cm}

	\colorbox{black!5}{%
		\parbox{0.75\textwidth}{%
			\vspace{0.4cm}
			\begin{tabular}{rl}
				\textbf{Alumno:}   & \parbox[t]{8cm}{\autores} \\[0.6cm]
				\textbf{Profesor:} & \profesor                 \\[0.25cm]
				\textbf{Fecha:}    & \fecha                    \\
			\end{tabular}
			\vspace{0.4cm}
		}%
	}

	\vfill

	{\color{black!70}\rule{\textwidth}{1.5pt}}

	\vspace{0.3cm}

\end{center}

\newpage

\section{Descripción General del Proyecto}
KitchenFlow es un sistema integral de control de producción diseñado para la gestión operativa y financiera de una cafetería. El sistema permite transformar materia prima en productos terminados mediante un recetario detallado con cálculo automático de costos y márgenes de utilidad.

\section{Usuarios y Seguridad}
El sistema implementa un esquema de seguridad basado en roles (RBAC) gestionado a través de \textbf{JWT (JSON Web Tokens)} y políticas de acceso a nivel de base de datos con \textbf{ZenStack}.

\begin{itemize}
	\item \textbf{Administrador:} Acceso total al sistema, gestión de usuarios, costos base y analítica financiera (\textit{véase fig. \ref{fig:seq_auth}}).
	\item \textbf{Personal de Cocina:} Registra producción y mermas; puede consultar recetas pero no modificar precios base.
	\item \textbf{Personal de Piso (Mesero):} Interfaz restringida únicamente al despacho de productos terminados.
\end{itemize}

\begin{figure}[H]
	\centering
	\fcolorbox{black}{white}{\includegraphics[height=0.3\textheight,keepaspectratio]{diagramas/auth_rbac_flow}}
	\caption{Secuencia de Autenticación y Autorización basada en Roles.}
	\label{fig:seq_auth}
\end{figure}

\section{Arquitectura y Stack Tecnológico}
Se utiliza una arquitectura desacoplada con \textbf{Hono} (Backend) y \textbf{Angular} (Frontend). La integridad de los tipos se garantiza compartiendo esquemas de validación entre ambas capas, asegurando que cualquier cambio en la estructura de datos sea detectado en tiempo de compilación.

El diagrama de la Fig. \ref{fig:arquitectura} resume tres capas coordinadas:
\begin{itemize}
	\item \textbf{Cliente Angular 18}: Componentes UI y señales que consumen contratos RPC generados a partir de esquemas Zod compartidos.
	\item \textbf{Servicios Hono}: Middleware JWT, controladores y contratos tipados bajo el mismo paquete de DTOs, evitando divergencias entre frontend y backend.
	\item \textbf{Persistencia ZenStack/Prisma}: Políticas declarativas que se ejecutan antes de tocar PostgreSQL 16, manteniendo coherencia de negocio en cada transacción.
\end{itemize}

\begin{figure}[H]
	\centering
	\fcolorbox{black}{white}{\includegraphics[height=0.5\textheight,keepaspectratio]{diagramas/Arquitectura}}
	\caption{Arquitectura del sistema KitchenFlow destacando el flujo de datos tipados.}
	\label{fig:arquitectura}
\end{figure}

\section{Especificación de Requerimientos Funcionales}

\subsection{Recetario y Gestión de Costos Promedio}
A diferencia de sistemas estáticos, KitchenFlow utiliza la metodología de \textbf{Costo Promedio Ponderado} para el inventario de insumos. Si el precio de un ingrediente varía entre compras, el sistema recalcula el costo unitario de la receta automáticamente (\textit{véase fig. \ref{fig:seq_recetario}}).
\begin{itemize}
	\item \textbf{Conversión de Unidades:} Manejo de factores de conversión entre unidades de compra (ej. Bulto 20kg) y consumo (ej. gramos).
\end{itemize}

\begin{figure}[H]
	\centering
	\fcolorbox{black}{white}{\includegraphics[height=0.3\textheight,keepaspectratio]{diagramas/recipe_cost_editor}}
	\caption{Secuencia de Configuración de Recetas y Costeo.}
	\label{fig:seq_recetario}
\end{figure}

\subsection{Motor de Explosión de Materiales y Producción}
El registro de producción activa una transacción atómica que valida el stock físico antes de realizar la "explosión" de la receta, transformando los insumos en existencias de vitrina (\textit{véase fig. \ref{fig:seq_produccion}}).

\begin{figure}[H]
	\centering
	\fcolorbox{black}{white}{\includegraphics[height=0.3\textheight,keepaspectratio]{diagramas/production_batch}}
	\caption{Secuencia de Registro de Producción (Explosión de Materiales con transacción).}
	\label{fig:seq_produccion}
\end{figure}

\subsection{Control de Mermas y Despacho}
El sistema permite identificar fugas mediante el cruce entre la producción registrada y las salidas marcadas por el personal de piso (\textit{véase fig. \ref{fig:seq_ventas}} y \ref{fig:seq_mermas}).

\begin{figure}[H]
	\centering
	\fcolorbox{black}{white}{\includegraphics[height=0.3\textheight,keepaspectratio]{diagramas/product_stock_sale}}
	\caption{Secuencia de Despacho de Ventas (ajuste inmediato de stock de vitrina).}
	\label{fig:seq_ventas}
\end{figure}

\begin{figure}[H]
	\centering
	\fcolorbox{black}{white}{\includegraphics[height=0.3\textheight,keepaspectratio]{diagramas/waste_stock_adjustment}}
	\caption{Secuencia de Registro de Mermas.}
	\label{fig:seq_mermas}
\end{figure}

\subsection{Abastecimiento y Costo Promedio Ponderado}
Cada registro de compra actualiza el costo promedio del ingrediente antes de validarse una nueva producción. El flujo de la Fig. \ref{fig:seq_abastecimiento} detalla cómo se recalcula el WAC y se actualiza ZenStack/Prisma.

\begin{figure}[H]
	\centering
	\fcolorbox{black}{white}{\includegraphics[height=0.3\textheight,width=0.85\textwidth,keepaspectratio]{diagramas/purchase_cost_recalculation}}
	\caption{Secuencia de Abastecimiento de Materia Prima.}
	\label{fig:seq_abastecimiento}
\end{figure}

\subsection{Dashboard de Rentabilidad}
La analítica mensual agrega datos de ventas, producción y mermas para mostrar KPIs accionables (\textit{véase fig. \ref{fig:seq_dashboard}}).

\begin{figure}[H]
	\centering
	\fcolorbox{black}{white}{\includegraphics[height=0.3\textheight,keepaspectratio]{diagramas/monthly_profitability}}
	\caption{Secuencia de Consulta del Dashboard de Rentabilidad.}
	\label{fig:seq_dashboard}
\end{figure}

\section{Modelo de Datos}
El esquema se define mediante ZenStack, integrando las políticas de seguridad directamente en las entidades de \texttt{User}, \texttt{Ingredient}, \texttt{Recipe} y \texttt{ProductionBatch}. La Fig. \ref{fig:clases} detalla las relaciones necesarias para costeo promedio, explosión de materiales y trazabilidad de compras/mermas.

Los puntos clave del modelo son:
\begin{itemize}
	\item \textbf{RecipeItem} actúa como tabla puente para registrar cantidades exactas por insumo dentro de cada receta.
	\item \textbf{ProductionBatch} conserva un \textit{snapshot} del costo calculado en el momento de producir, permitiendo auditorías posteriores.
	\item \textbf{PurchaseRecord} y \textbf{Waste} alimentan los ajustes de inventario y el cálculo del Costo Promedio Ponderado.
\end{itemize}

\begin{figure}[H]
	\centering
	\fcolorbox{black}{white}{\includegraphics[height=0.3\textheight,keepaspectratio]{diagramas/Clases}}
	\caption{Diagrama de Clases extendido con gestión de usuarios y roles.}
	\label{fig:clases}
\end{figure}
\section{Bocetos de Interfaces (Mockups)}

\subsection{[mp-01] Dashboard de Rentabilidad}
Pantalla de control para usuarios ADMIN con KPIs, comparativas de Ventas vs Producción vs Mermas y tabla de alertas.

\begin{figure}[H]
	\centering
	\fcolorbox{black}{white}{\includegraphics[width=0.6\textwidth,keepaspectratio]{mockups/Dashboard}}
	\\[0.3cm]
	\fcolorbox{black}{white}{\includegraphics[width=0.6\textwidth,keepaspectratio]{mockups/Dashboard2}}
	\caption{Dashboard mensual con KPIs y gráficos detallados (mockup combinado).}
	\label{fig:mock_dashboard}
\end{figure}

\subsection{[mp-02] Editor de Recetas y Costeo}
Formulario maestro-detalle para definir fichas técnicas, márgenes en tiempo real y tabla de insumos.

\begin{figure}[H]
	\centering
	\fcolorbox{black}{white}{\includegraphics[width=0.75\textwidth,keepaspectratio]{mockups/Recetas}}
	\caption{Editor de Recetas y Costeo con cálculo dinámico de margen.}
	\label{fig:mock_recetas}
\end{figure}

\subsection{[mp-03] Interfaz de Producción y Mermas}
Flujo unificado para programar lotes y registrar mermas con timeline auditado.

\begin{figure}[H]
	\centering
	\fcolorbox{black}{white}{\includegraphics[width=0.75\textwidth,keepaspectratio]{mockups/Mermas_y_produccion}}
	\caption{Panel dual de Producción y Mermas con acciones rápidas.}
	\label{fig:mock_mermas}
\end{figure}

\subsection{Panel de Ventas en Piso}
Terminal táctil para meseros con mosaico de productos y alertas de stock.

\begin{figure}[H]
	\centering
	\fcolorbox{black}{white}{\includegraphics[width=0.75\textwidth,keepaspectratio]{mockups/ventas}}
	\caption{Panel de Ventas en Piso con decremento inmediato de stock.}
	\label{fig:mock_ventas}
\end{figure}

\subsection{Terminal de Abastecimiento}
Wizard de recepción para recalcular costo promedio ponderado y documentar facturas.

\begin{figure}[H]
	\centering
	\fcolorbox{black}{white}{\includegraphics[width=0.75\textwidth,keepaspectratio]{mockups/Abastecimiento1}}
	\caption{Abastecimiento - Paso 1: datos de factura y proveedor.}
	\label{fig:mock_abast1}
\end{figure}

\begin{figure}[H]
	\centering
	\fcolorbox{black}{white}{\includegraphics[width=0.75\textwidth,keepaspectratio]{mockups/Abastecimiento2}}
	\caption{Abastecimiento - Paso 2: selección de insumos y cantidades.}
	\label{fig:mock_abast2}
\end{figure}

\begin{figure}[H]
	\centering
	\fcolorbox{black}{white}{\includegraphics[width=0.75\textwidth,keepaspectratio]{mockups/Abastecimiento3}}
	\caption{Abastecimiento - Paso 3: revisión de costos y cálculo WAC.}
	\label{fig:mock_abast3}
\end{figure}

\begin{figure}[H]
	\centering
	\fcolorbox{black}{white}{\includegraphics[width=0.75\textwidth,keepaspectratio]{mockups/Abastecimiento4}}
	\caption{Abastecimiento - Paso 4: confirmación y alerta de stock actualizado.}
	\label{fig:mock_abast4}
\end{figure}

\end{document}

% =====================================================
% ANEXO: Diferencias entre la propuesta y la implementación
% =====================================================
\clearpage
\section*{Anexo: Diferencias entre la propuesta y la implementación}
\addcontentsline{toc}{section}{Anexo: Diferencias entre la propuesta y la implementación}

\subsection*{Resumen ejecutivo}
En esta sección se documentan las desviaciones relevantes entre la propuesta original y la
implementaci\'on actual del proyecto KitchenFlow, la justificaci\'on de cada cambio y las
recomendaciones pendientes para alinear completamente la implementaci\'on con la propuesta.

\subsection*{Cambios principales}
\begin{itemize}
  \item \textbf{Stack y Arquitectura:} Propuesta original: Hono + ZenStack/Prisma + PostgreSQL.
	  Implementaci\'on: Express + Mongoose + MongoDB + OpenAPI. Justificaci\'on: se prioriz\'o
	  rapidez de desarrollo y despliegue con tecnolog\'ias ya disponibles en el equipo. Se
	  mantiene integridad de tipos mediante OpenAPI en lugar de Zod/ZenStack.

  \item \textbf{Autenticaci\'on y RBAC:} La propuesta planteaba pol\'iticas declarativas (ZenStack).
	  En la implementaci\'on se usa JWT firmado con helper propio y verificaciones en middleware
	  y controladores. Estado: funcional (roles ADMIN/KITCHEN/FLOOR presentes).

  \item \textbf{Costeo (WAC) y Abastecimiento:} El c\'alculo de Costo Promedio Ponderado (WAC)
	  est\'a implementado en el endpoint de compras (`/api/purchase-records`) y actualiza
	  `averageCost` y stock del insumo. Falta: factor de conversi\'on de unidades (p.ej. kg→g).

  \item \textbf{Explosi\'on de materiales y Producci\'on:} El flujo de producci\'on (reserva,
	  inicio, completado, cancelaci\'on) est\'a implementado en `production-batches`. Sin
	  embargo, no se usan transacciones at\'omicas de BD (riesgo de condiciones de carrera en
	  concurrencia alta).

  \item \textbf{Mermas y Ventas:} Registro de mermas y decremento de stock por ventas implementado
	  y visible en el dashboard anal\'itico.

  \item \textbf{Contratos y Cliente:} En lugar de compartir esquemas Zod, la integraci\'on se hizo
	  mediante OpenAPI y cliente generado (`frontend/src/app/api/schema.d.ts` y `client.ts`).

  \item \textbf{Mockups y Diagramas:} Todos los activos y el template LaTeX de la propuesta se
	  conservaron y se copiaron a `report/propuesta` para generar el reporte entregable.
\end{itemize}

\subsection*{Elementos no implementados o parciales}
\begin{itemize}
  \item \textbf{Transacciones at\'omicas en operaciones de producci\'on/compra:} Parcial — las
	  operaciones realizan varias escrituras por documento pero no usan sesiones/transactions de
	  MongoDB; se recomienda activarlas si se espera concurrencia.

  \item \textbf{Conversi\'on de unidades:} No implementada — si la operaci\'on maneja unidades de
	  compra distintas a las de consumo (bultos vs gramos) conviene añadir campos y reglas de
	  conversi\'on.

  \item \textbf{Pol\'iticas declarativas (ZenStack):} No implementadas — las reglas de negocio y
	  control de acceso est\'an codificadas manualmente.

  \item \textbf{Pruebas automatizadas:} Faltan tests unitarios e integraci\'on cr\'iticos.
\end{itemize}

\subsection*{Recomendaciones}
\begin{enumerate}
  \item Implementar transacciones de MongoDB en los endpoints cr\'iticos (`/api/production-batches`,
	  `/api/purchase-records`) para garantizar atomicidad en la reserva y consumos de stock.
  \item Añadir soporte de conversi\'on de unidades en `Ingredient`/`Recipe` y en el editor de
	  recetas para evitar errores de costeo cuando las compras vienen en unidades diferentes.
  \item Añadir pruebas autom\'aticas (tests) que cubran login, WAC, creaci\'on/completado de lote
	  y flujo de ventas.
  \item Documentar en el README la desviaci\'on de stack (por qu\'e se cambi\'o a Mongo/Express)
	  y los pasos para migrar a la stack original si fuera necesario.
\end{enumerate}

\subsection*{Estado del reporte}
El repositorio contiene ahora una copia completa de la propuesta en `report/propuesta` incluyendo
el PDF original y los activos; el present e archivo LaTeX se actualiz\'o con este anexo para que el
reporte compile con la informaci\'on de desviaciones y recomendaciones.

